% Part: many-valued-logic
% Chapter: three-valued-logics
% Section: lukasiewicz

\documentclass[../../../include/open-logic-section]{subfiles}

\begin{document}

\olfileid[pt]{mvl}{thr}{luk}

\olsection{Lógica de \L ukasiewicz}

Uma das primeiras propostas publicadas e desenvolvidas de uma lógica
multivalorada deve-se ao filósofo polonês Jan \L ukasiewicz, em 1921. \L
ukasiewicz foi motivado pelo problema aristotélico da batalha naval: parece
que, \emph{hoje}, a sentença ``Haverá uma batalha naval
amanhã'' não é verdadeira nem falsa: seu valor de verdade ainda não está
determinado. \L ukasiewicz propôs introduzir um terceiro valor de verdade
para tais sentenças ``contingentes futuras''.
\begin{quote}
Posso supor, sem contradição, que minha presença em Varsóvia em um
determinado momento do próximo ano, p.ex., ao meio-dia de 21 de dezembro, não está,
no presente, determinada nem positiva nem negativamente. Portanto, é
possível, mas não necessário, que eu esteja em Varsóvia no momento
indicado. Sob essa suposição, a proposição ``Estarei em Varsóvia
ao meio-dia de 21 de dezembro do próximo ano'' não pode, no presente,
ser verdadeira nem falsa. Pois, se fosse verdadeira agora, minha presença futura em
Varsóvia teria de ser necessária, o que contradiz a
suposição. Se fosse falsa agora, por outro lado, minha presença futura
em Varsóvia teria de ser impossível, o que também
contradiz a suposição. Portanto, a proposição considerada
não é, neste momento, verdadeira nem falsa e deve possuir um terceiro
valor, diferente de ``0'', ou falsidade, e de ``1'', ou verdade. Podemos
designar esse valor por ``$\frac{1}{2}$''. Ele representa ``o possível''
e se junta a ``o verdadeiro'' e ``o falso'' como terceiro valor.
\end{quote}
Usaremos $\Undef$ para o terceiro valor de verdade de \L
ukasiewicz.\footnote{\L ukasiewicz usa aqui ``possível'' de um modo
incomum hoje, a saber, no sentido de possível, mas não necessário.}

As funções de verdade dos conectivos $\lnot$, $\land$ e $\lor$
são fáceis de determinar segundo essa interpretação: a negação de uma sentença
contingente futura também é uma sentença contingente futura; portanto,
$\tf{\lnot}(\Undef) = \Undef$. Se um dos membros de uma
conjunção é indeterminado e o outro é verdadeiro, a conjunção também é
indeterminada---afinal, conforme o resultado futuro do membro
contingente, a conjunção pode acabar sendo verdadeira ou
falsa. Portanto, 
\[
    \tf{\land}(\True, \Undef) =
\tf{\land}(\Undef, \True) =
\Undef.
\]
Se o outro membro for falso, contudo, a conjunção não poderá
vir a ser verdadeira; portanto, \[\tf{\land}(\False, \Undef) =
\tf{\land}(\False, \Undef) = \False.\]
Os demais valores (quando os argumentos são valores de verdade determinados, $\True$
ou $\False$, são os mesmos da lógica clássica.

Para o condicional, a situação é um pouco mais delicada. Suponha que $q$
seja uma sentença contingente futura. Se $p$ for falsa, então $p \lif q$
será verdadeira, seja qual for o resultado de $q$; portanto, devemos definir
$\tf{\lif}(\False, \Undef) = \True$. E, se $p$ for verdadeira, então
$q \lif p$ será verdadeira, seja qual for o resultado de $q$; portanto,
$\tf{\lif}(\Undef, \True) = \True$. Se $p$ for verdadeira, então $p
\lif q$ poderá acabar sendo verdadeira ou falsa; logo, $\tf{\lif}(\True,
\Undef) = \Undef$. De modo semelhante, se $p$ for falsa, então $q
\lif p$ poderá acabar sendo verdadeira ou falsa; logo,
$\tf{\lif}(\Undef, \False) = \Undef$. Resta o
caso em que $p$ e $q$ são ambas contingentes futuras. Com base na
motivação, deveríamos realmente atribuir $\Undef$ nesse caso.
Entretanto, isso faria com que $!A \lif !A$ \emph{não} fosse uma tautologia. \L
ukasiewicz não teve dificuldade em abandonar $!A \lor \lnot !A$ e $\lnot(!A
\land \lnot !A)$, mas hesitou em abandonar $!A \lif !A$. Por isso,
estipulou $\tf{\lif}(\Undef, \Undef) =
\True$.

\begin{defn}\ollabel{def:lukasiewicz}
A lógica trivalorada de \L ukasiewicz é definida pela matriz:
\begin{enumerate}
  \item A linguagem proposicional padrão $\Lang L_0$, com
  $\lnot$, $\land$, $\lor$, $\lif$.
  \item O conjunto de valores de verdade $V = \{\True, \Undef, \False\}$.
  \item $\True$ é o único valor designado, O.s.\@, $V^+ = \{\True\}$.
  \item As funções de verdade são dadas pelas seguintes tabelas:
  \begin{center}
    \begin{tabular}{c|c} 
      $\tf{\lnot}$ & \\ 
      \hline  
      $\True$ & $\False$ \\ 
      $\Undef$ & $\Undef$ \\
      $\False$ & $\True$ 
    \end{tabular}
    \quad
    \begin{tabular}{c|ccc} 
      $\tf{\land}[\LogLuk[3]]$ & $\True$ & $\Undef$ & $\False$ \\ 
      \hline 
      $\True$ & $\True$ & $\Undef$ & $\False$ \\ 
      $\Undef$ & $\Undef$ & $\Undef$ & $\False$\\ 
      $\False$ & $\False$ & $\False$ & $\False$ 
    \end{tabular}
    \\[2ex]
    \begin{tabular}{c|ccc} 
      $\tf{\lor}[\LogLuk[3]]$ & $\True$ & $\Undef$ & $\False$ \\ 
      \hline 
      $\True$ & $\True$ & $\True$ & $\True$ \\ 
      $\Undef$ & $\True$ & $\Undef$ & $\Undef$ \\
      $\False$ & $\True$ & $\Undef$ & $\False$ 
    \end{tabular}
    \quad
    \begin{tabular}{c|ccc} 
      $\tf{\lif}[\LogLuk[3]]$ & $\True$ & $\Undef$ & $\False$ \\ 
      \hline 
      $\True$ & $\True$ & $\Undef$ & $\False$ \\ 
      $\Undef$ & $\True$ & $\True$ & $\Undef$  \\ 
      $\False$ & $\True$ & $\True$ & $\True$ 
    \end{tabular}
  \end{center} 
\end{enumerate}
\end{defn}

Como se vê facilmente, qualquer !!{formula}~$!A$ que contenha apenas $\lnot$,
$\land$ e $\lor$ terá o valor de verdade~$\Undef$ se todas
as suas !!{propositional variable}s receberem~$\Undef$. Assim, por
exemplo, as tautologias clássicas $p \lor \lnot p$ e $\lnot(p
\land \lnot p)$ não são tautologias em $\LogLuk[3]$, pois $\pValue{v}(!A)
= \Undef$ sempre que $\pAssign v(p) = \Undef$.

Nas !!{valuation}s em que $\pAssign v(p) = \True$ ou $\False$, $\pValue
v(!A)$ coincidirá com seu valor de verdade clássico.

\begin{prop}
  Se $\pAssign v(p) \in \{\True, \False\}$ para todo $p$ em~$!A$, então
  $\pValue v(!A)[\LogLuk[3]] = \pValue v(!A)[\LogCL]$.
\end{prop}

\begin{prob}\label{mvl:thr:luk:prob:luk-iff} Suponha que definamos
  $\pValue v(!A \liff !B) = \pValue v((!A \lif !B) \land (!B \lif
  !A))$ em~$\LogLuk[3]$. Qual seria a tabela-verdade de $\liff$?
\end{prob}

Muitas tautologias clássicas \emph{são} também tautologias em \LogLuk[3],
por exemplo, $\lnot p \lif (p \lif q)$. Assim como na lógica clássica, podemos
usar tabelas-verdade para verificar isso:
\begin{center}
\begin{tabular}{cc|cccccc}
  $p$ & $q$ & $\lnot$ & $p$ & $\lif$ & $\smash{(}p$ & $\lif$ & $q\smash{)}$ \\ \hline
  \True & \True & \False & \True & \True & \True & \True & \True \\
  \True & \Undef & \False & \True & \True & \True & \Undef & \Undef \\
  \True & \False & \False & \True & \True & \True & \False & \False \\
  \Undef & \True & \Undef & \Undef & \True & \Undef & \True & \True \\
  \Undef & \Undef & \Undef & \Undef & \True & \Undef & \True & \Undef \\
  \Undef & \False & \Undef & \Undef & \True & \Undef & \Undef & \False \\
  \False & \True & \True & \False & \True & \False & \True & \True \\
  \False & \Undef & \True & \False & \True & \False & \True & \Undef \\
  \False & \False & \True & \False & \True & \False & \True & \False \\  
\end{tabular}
\end{center}

\begin{prob}
  Mostre que as seguintes fórmulas são tautologias em \LogLuk[3]:
  \begin{enumerate}
    \item $p \lif (q \lif p)$
    \item\label{mvl:thr:luk:prob:luk-taut-2} $\lnot(p \land q) \liff (\lnot p \lor \lnot q)$
    \item\label{mvl:thr:luk:prob:luk-taut-3} $\lnot(p \lor q) \liff (\lnot p \land \lnot q)$
  \end{enumerate}
  (Em \olref[mvl][thr][luk]{prob:luk-taut-2} e
  \olref[mvl][thr][luk]{prob:luk-taut-3}, tome $!A \liff !B$ como uma
  abreviação de $(!A \lif !B) \land (!B \lif !A)$, ou consulte sua
  solução de \olref[mvl][thr][luk]{prob:luk-iff}.)
\end{prob}

\begin{prob}
  Mostre que as seguintes tautologias clássicas não são tautologias em~\LogLuk[3]:
  \begin{enumerate}
    \item $(\lnot p \land p) \lif q)$
    \item $((p \lif q) \lif p) \lif p$
    \item $(p \lif (p \lif q)) \lif (p \lif q)$
  \end{enumerate}
\end{prob}

Poder-se-ia então pensar que, embora nem todas as tautologias clássicas
sejam tautologias em~$\LogLuk[3]$, elas deveriam ao menos assumir
o valor~$\True$ ou o valor~$\Undef$ em toda
!!{valuation}. Não é assim. Um contraexemplo é dado por
\[
  \lnot(p \lif \lnot p) \lor \lnot(\lnot p \lif p)
\]
que é $\False$ se $p$ for~$\Undef$.

\begin{prob}
  Quais das relações a seguir valem na lógica de \L ukasiewicz? Dê uma tabela-verdade para cada uma.
  \begin{enumerate}
    \item $p, p \lif q \Entails q$
    \item $\lnot\lnot p \Entails p$
    \item $p \land q \Entails p$
    \item $p \Entails p \land p$
    \item $p \Entails p \lor q$
  \end{enumerate}
\end{prob}

\L ukasiewicz esperava construir uma lógica da possibilidade com base em seu
sistema trivalorado, introduzindo um conectivo unário $\Diamond
!A$ (para ``$!A$ é possível'') e um $\Box !A$ correspondente (para ``$!A$
é necessário''):
\begin{center}
  \begin{tabular}{c|c} 
    $\tf{\Diamond}$ & \\ 
    \hline  
    $\True$ & $\True$ \\ 
    $\Undef$ & $\True$ \\
    $\False$ & $\False$ 
  \end{tabular}
  \quad
  \begin{tabular}{c|c} 
    $\tf{\Box}$ & \\ 
    \hline  
    $\True$ & $\True$ \\ 
    $\Undef$ & $\False$ \\
    $\False$ & $\False$ 
  \end{tabular}
\end{center}
Em outras palavras, $p$ é possível se, e somente se, ainda não estiver determinado como falso;
e $p$ é necessário se, e somente se, já estiver determinado como verdadeiro.

\begin{prob}
  Mostre que $\Box p \liff \lnot\Diamond \lnot p$ e $\Diamond p \liff
  \lnot \Box \lnot p$ são tautologias em~\LogLuk[3], estendida com as
  tabelas-verdade de $\Box$ e~$\Diamond$.
\end{prob}

Contudo, as deficiências dessa lógica modal proposta logo se tornaram
evidentes: qualquer que seja o resultado, $p \land \lnot p$ nunca poderá
ser verdadeira. Portanto, mesmo que isso não esteja agora determinado (e seja, por isso,
indeterminado), deve contar como impossível, O.s.\@, $\lnot \Diamond(p
\land \lnot p)$ deveria ser uma tautologia. No entanto, se $\pAssign v(p) =
\Undef$, então $\pValue v(\lnot \Diamond(p \land \lnot p)) =
\Undef$. Embora \L ukasiewicz estivesse certo ao afirmar que dois valores de
verdade não bastariam para acomodar distinções modais como
possibilidade e necessidade, introduzir um terceiro valor de verdade tampouco
basta.

\end{document}
